\section{DyanamoDB}

\subsection{DynamoDB - in short}
\begin{itemize}
    \item NoSQL database, fully managed, massive scale (1,000,000 requests per second)
    \item Similar to Apache Cassandra (can migrate to DynamoDB)
    \item No disk space to provision, max object size is 400kb
    \item Capacity: provisioned (WCU, RCU and Auto Scaling) or on-demand
    \item Support CRUD (Create, Read, Update, Delete)
    \item Read: Eventually or strong consistency
    \item Supports transaction across multiple tables (ACID support)
    \item Backups available, point in time recovery
    \item Table classes: Standard and Infrequent Access (IA)
\end{itemize}

\subsection{DynamoDB - Basics}
\begin{itemize}
    \item DynamoDB is made of tables
    \item Each table has a primary key (must be decided at creation time)
    \item Each table can have aan infinite number of times (= rows)
    \item Each item has attributes (can be added over time - can be null)
    \item Maximum size of a item is a 400kb
    \item Data types supported are:
    \begin{itemize}
        \item ScalarTypes: String, Number, Binary, Boolean, Null
        \item Document Types: List, Map
        \item Set Types: String Set, Number Set, Binary Set
    \end{itemize}
\end{itemize}

\subsection{DynamoDB - Primary Keys}
%TODO add diagram
\begin{itemize}
    \item Option 1: Partition Key only (HASH)
    \item Partition key must be unique for each item
    \item Partition key must be "diverse" so that the data is distributed
    \item Example: user_id for a users table
\end{itemize}


%TODO add diagram
\begin{itemize}
    \item Option 2: Partition Key + Sort Key
    \item The combination must be unique
    \item Data is grouped by partition key
    \item Sort key == range key
    \item Example: user-games table
    \item user_id for the partition key
    \item game_id for the sort key
\end{itemize}

\item \textbf{Example good sort key: timestamp}

\subsection{DynamoDB - Indexes}
\begin{itemize}
    \item Object = partitioon key + optional sort key + attributes
    \item LSI - Local Secondary Index
    \item Keep the same parition key
    \item Select an alternative key
    \item Must be defined at table creation time
    \item GSI - Global Secondary Index
    \item Change the partition key and optionally the sort key
    \item Can be defined after the table is created
    \item \textbf{You can only query by PK + sort key on the main table and indexes \neq RDS}
\end{itemize}

\subsection{Amazon Kinesis Data Streams for DynamoDB}
%TODO add diagram
\begin{itemize}
    \item You can use Kinesis Data Streams to capture item-level changes in DynamoDB
    \item Customer and longer data retention period (> 24 hours in DynamoDB Streams)
\end{itemize}

\subsection{DynamoDB Solution Architecture Indexing Objects in DynamoDB}
%TODO add diagram
\begin{itemize}
    API for object metadata
    \item Search by date
    \item Total storage used by a customer
    \item List of all object with certain attributes
    \item Find all objects uploaded within a date range
\end{itemize}

\subsection{DynamoDB - DAX}
%TODO add diagram
\begin{itemize}
    \item DAX = DynamoDB Accelerator
    \item Seamless cache for DynamoDB, no application re-write
    \item Writes go through DAX to DynamoDB
    \item Micro second latency for cached reads and queries
    \item Solives the Hot Key problem (too many reads)
    \item 5 minutes TTL for cache by default
    \item Up to 10 nodes in the cluster
    \item Multi AZ (3 nodes minimuum reccommended for production)
    \item Secure (Encryption at rest with KMS, VPC, IAM, CloudTrail....)
\end{itemize}

\subsection{DynamoDB - DAX vs ElastiCache}
%TODO add diagram


\section{Amazon OpenSearch}

\subsection{Amazon OpenSearch (ex ElasticSearch)}
\begin{itemize}
    \item New name is Amazon OpenSearch
    \item ElasticSearch => OpenSearch
    \item Kibana => OpenSearch Dashboards
    \item Managed versions of OpenSearch (open-source project, fork of ElastiSearch)
    \item Two modes: Managed cluster or Serverless cluster
    \item Use cases:
    \item Log Analytics
    \item Real time application monitoring
    \item Security Analytics
    \item Full Text Search
    \item Clickstream Analytics
    \item Indexing
\end{itemize}

\subsection{OpenSearch + OS Dashboards + Logstash}
\begin{itemize}
    \item OpenSearch (ex ElasticSearch): Provide search and indexing capability
    \item OpenSearch Dashboard (ex Kibana)
    \item Provide real-time dashboards on top of the data that sits in OpenSearch
    \item Alternative to CloudWatch dashboards (more advanced capabilities)
    \item Logstash:
    \item Log ingestion mechanism, use the "Logstash Agent"
    \item Alternative to CloudWatch Log (you decide on retention and granularity)
\end{itemize}

\subsection{OpenSearch patterns DynamoDB}
%TODO see diagram

\subsection{OpenSearch patterns CloudWatch Logs}
%TODO see diagram


\section{RDS}
\begin{itemize}
    \item Engines: PostgreSQL, MySQL, MariaDB, IBM DB2, Oracle, SQL Server
    \item Managed DB: provisioning, backups, patching, monitoring
    \item Launched within a VPC, usually in a private subnet, control network access using security groups (important when using Lambda)
    \item Storage by EBS, can increase volume size with auto-scaling
    \item Backups: automated with point-in-time recovery. Backups expire
    \item Snapshots: manual, can make copies of snapshots cross region
    \item RDS Events: get notified via SNS for events (operations, outages.....)
\end{itemize}

\subsection{RDS - Multi AZ and Read Replicas}

\subsubsection{Multi-AZ}
%TODO add diagram
- Standby instance for failover in case of outage

\subsubsection{Read Replicas}
%TODO add diagram
- Increase read, throughput. Eventual consistency. Can be cross-region

\subsection{RDS - Distributing Reads across Replicas}
%TODO add/see diagram

\subsection{RDS - Security (reminder)}
\begin{itemize}
    \item KMS encryption at rest for underlying EBS volumes / snapshots
    \item Transparent Data Encryption (TDE) for Oracle and SQL Server
    \item SSL encryption to RDS is possible for all DB (in - flight)
    \item IAM authentication for MySQL, PostgreSQL and MariaDB
    \item Authorisation still happens within RDS (not in IAM)
    \item Can copy an un-encrypted RDS snapshot into an encrypted one
    \item CloudTrail cannot be used to track queries made within RDS
\end{itemize}

\subsection{RDS - IAM Authetication}
%TODO add diagram
\begin{itemize}
    \item IAM database authentication works with MariaDB, MySQL and PostgreSQL
    \item You don't need a password, just an authentication token obtained through IAM and RDS API calls
    \item Auth token has a lifetime of 15 minutes
    \item Benefits
    \item Network in/out must be encrypted using SSL
    \item IAM to centrally manage users instead of DB
    \item Can leverage IAM Roles and EC2 instance profiles for easy integration
\end{itemize}

\subsection{About RDS for Oracle - Exam Tips}
%TODO see/add diagrams
\begin{itemize}
    \item Backups
    \item Use RDS backups for backups and restore to Amazon RDS for Oracle
    \item Use Oracle RMAN (Recovery Manager) for backups and restore to-non RDS (RDS not supported)
    \item Real Application Clusters (RAC)
    \item RDS for Oracle does NOT support RAC
    \item RAC is working on Oracle on EC2 instances because you have full control
    \item RDS for Oracle supports Transparent Data Encryption (TDE) to encrypt data before it's written to storage
    \item DMS work on Oracle RDS
\end{itemize}

\subsection{About RDS for MySQL}
%TODO add diagram
\begin{itemize}
    \item You can use the native mysqldump to migrate a MySQL RDS DB to non-RDS
    \item The external MySQL database can run either on-premises in your data center, or on an Amazon EC2 instance
\end{itemize}

\subsection{RDS Proxy for AWS Lambda}
%TODO add diagram
\begin{itemize}
    \item When using Lambda functions with RDS, it opens and maintains a database connection
    \item This can result in a "TooManyConnections" exception
    \item With RDS Proxy, you no longer need code that handles cleaning up idle connections and managing connection pools
    \item Supports IAM authentication or DB authentication, auto-scaling
    \item The lambda function must be deployed in you VPC, because RDS Proxy is never publicly accessible
\end{itemize}

\subsection{RDS Solution Architecture Cross Region Failover}
%TODO add/see diagram


\section{Aurora - Part 1}
\begin{itemize}
    \item DB Engines: PostgreSQL-compatible and MySQL- compatible
    \item Storage: Automatically grows up to 256TB, 6 copies of data, multi-AZ
    \item Read Replicas: up to 15 RR, reader endpoint to access them all
    \item Cross Region RR: entire database is copied (not select tables)
    \item Load / Offload data directly from / to S3: efficient uuse of resources
    \item Backup, Snapshots and Restore: same as RDS
\end{itemize}

\subsection{Aurora High Availability and Read Scaling}
%TODO add diagram
\begin{itemize}
    \item 6 copies of your data across 3 AZ
    \item 4 copies out of 6 needed for writess
    \item 3 copies out of 6 need for reads
    \item self healing with peer-to-peer replication
    \item storage is striped across 100s of volumes
    \item Automated failover for master for less that 30 seconds
    \item Master + up to 15 Aurora Read Replicas serve reads
    \item Support for Cross Region Replication
\end{itemize}

\subsection{Aurora DB Cluster}
%TODO add diagram

\subsection{Aurora Endpoints}
\begin{itemize}
    \item Endpoints = Host Address + Port
    \item Cluster Endpoint (Writer Endpoint)
    \item Connects to the current primary DB instance in the Aurora cluster
    \item Used for all write operations in the DB cluster (inserts, updates, deletes and queries)
    \item Reader Endpoint
    \item Provides load-balancing for read only connection to all Aurora Replicas in the Aurora cluster
    \item Used only for read operations (queries)
    \item Custom Endpoint
    \item Represent a set on DB instances that you choose in the Aurora cluster
    \item Used when you want to connect to different subset of DB instances with different capacities and configurations (e.g, different DB parameter group)
    \item Instance Endpoint
    \item Connects to specific DB instances in the Aurora cluster
    \item Used when you want to diagnosis and fine tune a specific DB instance
\end{itemize}

\subsection{Aurora - Customer Endpoints}
%TODO add diagram
\begin{itemize}
    \item Define a subset of Aurora Instances as a Customer Endpoint
    \item Example: Run analytical queries on specific replicas
    \item The reader endpoint is generally not used after defining Customer Endpoints
\end{itemize}

\subsection{Aurora Logs}
%TODO add diagram
\begin{itemize}
    \item You can monitor the following types of Aurora MySQL log files:
    \item Error log
    \item Slow query log
    \item General log
    \item The audit log
    \item These log files are either downloaded or published to CloudWatch Logs
\end{itemize}

\subsection{Troubleshooting RDS and Aurora Performance}
\begin{itemize}
    \item Performance Insights: find issues by waits, SQL statements, hosts and users
    \item CloudWatch Metrics: CPU, Memory, Swap Usage
    \item Enhanced Monitoring Metrics: at host level, process view, per second metric
    \item Slow Query Logs
\end{itemize}

\subsection{Aurora Serverless}
%TODO add diagram
\begin{itemize}
    \item Automated database instantiation and auto-scaling based on actual usage
    \item Good for infrequent, intermittent or unpredictable workloads
    \item No capacity planing needing
    \item Pay per second, can be most cost-effective
\end{itemize}

\subsection{Aurora Serverless - Data API}
%TODO add diagram
\begin{itemize}
    \item Access Aurora Serverless DB with a simple API endpoint (no JDBC connection needed)
    \item Secure HTTPS endpoint to run SQL statements
    \item No persistent DB connections management
    \item Users must be granted permissions to Data API and Secrets Manager (where credentials are checked)
\end{itemize}

\subsection{RDS Proxy for Aurora}
%TODO add diagram/diagrams
- Ability to create an additional read\itemonly endpoint that connects to Aurora Read Replicas only

\subsection{Global Aurora}
%TODO add diagram
\begin{itemize}
    \item Aurora Cross Region Read Replicas
    \begin{itemize}
        \item Useful for disaster recovery
        \item Simple to put in place
    \end{itemize}

    \item Aurora Global Database (recommended)
    \begin{itemize}
        \item 1 Primary Region (read/write)
        \item Up to 10 secondary (read-only) regions; replication lag is less than 1 second
        \item Up to 16 Read Replicas per secondary region
        \item Helps decrease latency
        \item Promoting another region (for disaster recovery) has an RTO of less than 1 minute
        \item Ability to manage the RPO in Aurora PostgreSQL
    \end{itemize}
\end{itemize}

\subsection{Aurora Global - Write Forwarding}
%TODO add diagram
\begin{itemize}
    \item Enables Secondary DB Clusters to forward SQL statements that perform write operations to the Primary DB Cluster
    \item Data is always changed first on the Primary DB Cluster, then replicated to the Secondary DB Clusters
    \item Primary DB Cluster always has an up-to-date copy of all data
    \item Reduce the number of endpoints
\end{itemize}

\subsection{Covert RDS to Aurora}
%TODO add / see diagram